%% =====================================================================
%%  T-18-05  The Man in the Gray Coat
%%  -------------------------------------------------------------------
%%  One idea per file. Core is mandatory. Slots in canonical order.
%%  Max 700 words. No layout commands. No filler. New source material
%%  for this entry goes in sources/B9/T-18-05.tex -- never by
%%  rewriting this file (docs/RULES.md R0.1).
%% =====================================================================
%% @kind: topic
%% @id: T-18-05
%% @title: The Man in the Gray Coat
%% @subtitle: Sixty-five percent obeyed a plain order to the end --- and every predictor said otherwise
%% @chapter: c18
%% @order: 50
%% @status: stable
%% @sources: B9
%% @updated: 2026-09-20

\topic{T-18-05}{The Man in the Gray Coat}{Sixty-five percent obeyed a plain order to the end --- and every predictor said otherwise}

\begin{core}
B9's laboratory stripped coercion to its verbal minimum: a university hall, a learning experiment, a calm stranger in a gray technician's coat reading four scripted prods. Twenty-six of forty ordinary men --- 65 percent --- carried every shock to the 450-volt maximum, two steps past the switch marked Danger. Refusal cost nothing but the experiment. Psychologists asked to predict the outcome estimated 1.2 percent.
\end{core}
\begin{mechanism}
Nothing in the situation forced anyone; the engine was structure. Authority arrived entirely as symbol and setting --- the university's sponsorship, the worthy purpose, the impassive manner: the trappings of T-18-01 shown at their extreme, with no enforcement behind them. The request escalated in 15-volt steps, so no single step was ever the decision; by the victim's last protest (a wall-pounding at 300 volts, silence from 315) the subject had already answered twenty times. The exits were closed by script, in ascending order: \emph{Please continue}; \emph{The experiment requires that you continue}; \emph{It is absolutely essential that you continue}; and finally the bare denial of any exit. Objections were absorbed into the frame rather than answered by it: the shocks were ``painful but not dangerous''; ``whether the learner likes it or not,'' the run had to finish. Obligation came pre-loaded: the subject had been paid (and told the money was his no matter what), roles had been drawn by lot, the victim had volunteered first, and the subject had taken a 45-volt sample shock on his own wrist. B9 enumerates thirteen such features and locates the conflict in two dispositions colliding --- the reluctance to harm against the pull to obey perceived authority. Obedience was not ease: sweating, trembling and stuttering were routine, fourteen of forty laughed nervously, three went into full seizures. The body dissented while the hand complied.
\end{mechanism}
\begin{conditions}
Strongest with sponsorship, a closed setting, small increments and scripted prods; B9 notes that later runs which dissociated the study from the university weakened it. Defiance existed --- fourteen refused --- but never early: no subject stopped before the victim first protested at 300 volts. Payment was not the engine: forty-three unpaid undergraduates, run separately, obeyed at very similar rates.
\end{conditions}
\begin{application}
  \item Wrap the demand in an institution and a purpose larger than the person; let the frame, not you, carry the request.
  \item Escalate in steps no one would refuse individually, and never state the endpoint.
  \item Script the prods from mild to absolute, and meet objections by restating the frame --- it is safe; that is not the point.
  \item Collect a small sunk cost early (payment, a signature, a first yes) so quitting reads as loss.
  \item Keep the setting closed: no time to reflect, no peers to consult.
\end{application}
\begin{feedback}
  \item They comply while signalling distress --- \emph{I can't do this}, while doing it.
  \item Objections seek permission and, denied, dissolve; the frame has an answer for everything.
  \item \emph{Just to the next step} recurs; the endpoint is never said aloud.
  \item They ask you whether to stop. The transfer of the decision is the tell.
\end{feedback}
\begin{failure}
The seizures are the receipt: visible collapse can break the frame outright. And the prod sequence is the mechanism's confession --- by the fourth prod (\emph{you have no other choice}), obedience is visibly being manufactured, and the fourteen refusals clustered exactly there. Afterward the frame is seen through at leisure; B9's own procedure arranged a friendly reconciliation with the victim, and the practitioner's accounts come due in their own currency (T-27-01).
\end{failure}
\begin{countermeasures}
  \item Hear the fourth prod coming: whoever must say \emph{you have no other choice} has none. Real authority states consequences; it does not deny exits.
  \item Price the exit before assuming it is locked. In the study, defiance cost \$4.50 and a walk to the door.
  \item Convert the ladder to its endpoint --- \emph{where does step thirty land?} --- before step five.
  \item Decide in cold blood what you will not do; a frame left open will decide in hot blood (T-24-05).
  \item Tension mid-compliance --- the laughter, the sweat --- is the gut alarm under production conditions (T-23-04). Audit before the next step, not after the last.
\end{countermeasures}

\sources{B9}
\seesources
\seealso{T-18-01, T-18-02, T-24-02}
